\section{Measures of Complexity}
\label{sec:measures}

Once we have the $\epsilon$-machine, we can calculate intrinsic properties of the process that go beyond simple statistics.

\subsection{Definitions}
We care about two fundamental quantities: structure and randomness.

\begin{description}
    \item[Statistical Complexity ($C_\mu$)] \hfill \\
    How much memory does the process need to store?
    It is defined as the Shannon entropy of the causal states themselves.
    \begin{equation}
        C_\mu = H[\mathcal{S}] = - \sum_{\sigma \in \mathcal{S}} \pi(\sigma) \log_2 \pi(\sigma)
    \end{equation}
    where $\pi(\sigma)$ is the probability of being in state $\sigma$.

    \item[Entropy Rate ($h_\mu$)] \hfill \\
    How random is the next symbol, given the past?
    It is the average uncertainty of the next symbol, averaged over the causal states.
    \begin{equation}
        h_\mu = H[X|\mathcal{S}] = \sum_{\sigma \in \mathcal{S}} \pi(\sigma) H[X|\sigma]
    \end{equation}
\end{description}

\begin{sidebar}{Intuition: Stored vs. Generated Information}
These two measures give us a complete picture of the process's information budget:
\begin{enumerate}
    \item $C_\mu$ is the information \textbf{stored} from the past. It represents \textit{Structure}.
    \item $h_\mu$ is the new information \textbf{generated} at each step. It represents \textit{Randomness}.
\end{enumerate}

\par\vspace{2mm}\hrule\vspace{2mm}

\textbf{Examples:}
\begin{itemize}
    \item \textbf{Clock}: High $C_\mu$ (many states to track), Zero $h_\mu$ (perfectly predictable).
    \item \textbf{Fair Coin}: Zero $C_\mu$ (no memory needed), High $h_\mu$ (maximum surprise).
\end{itemize}
\end{sidebar}

\subsection{Calculation}

Let's compute these values for our Golden Mean process ($p=0.5$).

\begin{sidebar}{Derivation for Golden Mean}
    \textbf{1. Statistical Complexity ($C_\mu$)}

    The stationary distribution over states is derived from the transition matrix. For the Golden Mean ($p=0.5$):
    \[ \pi(A) = 2/3, \quad \pi(B) = 1/3 \]
    The entropy of this distribution is:
    \begin{align*}
    C_\mu &= - \sum \pi(\sigma) \log_2 \pi(\sigma) \\
          &= - (2/3 \log_2 2/3 + 1/3 \log_2 1/3) \\
          &\approx 0.918 \text{ bits}
    \end{align*}

    \textbf{2. Entropy Rate ($h_\mu$)}

    This is the average uncertainty of the next symbol given the current state.
    \begin{itemize}
        \item State $A$ emits $0$ or $1$ with equal probability: $H[X|A] = 1$ bit.
        \item State $B$ always emits $1$: $H[X|B] = 0$ bits.
    \end{itemize}
    Weighted by state probability:
    \begin{align*}
    h_\mu &= \sum \pi(\sigma) H[X|\sigma] \\
          &= (2/3 \times 1) + (1/3 \times 0) \\
          &= 0.667 \text{ bits}
    \end{align*}
\end{sidebar}

\subsection{Comparison of Processes}
Different processes occupy different places on the complexity-entropy plane.

\begin{emictable}[label=tab:measures]{Measures for Common Processes}
\begin{tabular}{lp{3.5cm}ll}
\toprule
Process & Description & $C_\mu$ & $h_\mu$ \\
\midrule
\textbf{Fair Coin} & No memory, pure randomness. & 0.00 & 1.00 \\
\textbf{Period-2} & Alternating 0101... Needs 1 bit of memory. & 1.00 & 0.00 \\
\textbf{Golden Mean} & No '00'. A mix of structure and choice. & 0.92 & 0.67 \\
\bottomrule
\end{tabular}
\end{emictable}
